\section{Comparação consolidada}
\label{sec:comparacao}

Esta seção apresenta as três fontes de números lado a lado, sempre
sob o mesmo protocolo \emph{paper-faithful} de \S{}\ref{sec:protocolo}:

\begin{enumerate}[leftmargin=*,itemsep=0pt]
  \item \textbf{Paper}: Tabelas 10/11/12/13 do paper Alsaedi 2020.
  \item \textbf{Nossa reprodução \emph{paper-faithful}}: oito
        \emph{baselines} via \texttt{run\_baselines\_paper.py}.
  \item \textbf{WiSARD-family}: cinco modelos sem peso, best F1
        sobre os \emph{grids} de \emph{sweep}.
\end{enumerate}

\noindent Lembrete: as comparações desta seção são informativas (não
tentativa de reprodução byte-a-byte), pelo motivo da divergência de
tamanho documentada em \S{}\ref{sec:dataset_mismatch}.

\noindent As subseções primeiro visualizam essas fontes
(\emph{heatmaps}, \emph{bar chart}, \emph{wins matrix} e \emph{boxplot}
agregado); em seguida medem os dois eixos de custo (tempo e memória)
sob um \emph{harness} único, traçam a fronteira de Pareto
F1 $\times$ custo e encerram com a tabela-síntese e sua leitura.

\subsection{Master heatmap: per-device binário}

A Figura~\ref{fig:master_perdevice} é o \emph{heatmap} principal do
relatório: sete bases $\times$ três grupos $\times$ até oito modelos, F1
\emph{weighted} (coluna do meio = reprodução \emph{paper-faithful};
direita = os cinco WiSARDs). Os três grupos concordam base a base, e a
família WiSARD ocupa o mesmo regime dos \emph{baselines} não-lineares:
lidera as bases estruturadas junto de kNN/RF/CART, satura nas de
\emph{leak} e fica a cerca de 1 pp de RF/CART na média (leitura
detalhada na \S{}\ref{sec:leitura}).

\begin{figure}[ht]
  \centering
  \includegraphics[width=\linewidth]{v06_heatmap_perdevice.png}
  \caption{Master heatmap: F1 \emph{weighted} médio para per-device
           binário. Três grupos: paper, nossa reprodução, WiSARD-family.}
  \label{fig:master_perdevice}
\end{figure}

\subsection{Combined: binário e multi-classe}

A Figura~\ref{fig:master_combined} reúne as duas tarefas \emph{combined}
(binário à esquerda, multi-classe à direita). O \emph{gap} para RF/CART,
pequeno no \emph{per-device}, alarga-se aqui: a família cobre
$\approx 0{,}72$ no binário e cai para $\approx 0{,}41$--$0{,}43$ no
multi-classe, à frente dos modelos lineares mas atrás das árvores
(\S{}\ref{sec:leitura}).

\begin{figure}[ht]
  \centering
  \begin{subfigure}[t]{0.48\linewidth}
    \includegraphics[width=\linewidth]{v06_heatmap_combined_binary.png}
    \caption{Combined binário (Tab.\ 12 do paper).}
  \end{subfigure}\hfill
  \begin{subfigure}[t]{0.48\linewidth}
    \includegraphics[width=\linewidth]{v06_heatmap_combined.png}
    \caption{Combined multi-classe (Tab.\ 13).}
  \end{subfigure}
  \caption{Master heatmaps das duas tarefas \emph{combined}:
           paper $\times$ ours $\times$ WiSARD-family.}
  \label{fig:master_combined}
\end{figure}

\subsection{Bar chart: campeão de cada grupo por base}

A Figura~\ref{fig:bar_best_perbase} compacta o resultado anterior em um
\emph{bar chart}: para cada base, a altura é o melhor F1 de cada grupo.
O campeão WiSARD iguala ou supera o campeão \emph{baseline} em cinco das
sete bases \emph{per-device} (incluindo \texttt{Modbus}), ficando atrás
apenas em \texttt{GPS\_Tracker} e \texttt{Weather}.

\begin{figure}[ht]
  \centering
  \includegraphics[width=0.85\linewidth]{v06_bar_best_perbase.png}
  \caption{Best F1 por base, mostrando o campeão de cada um dos três
           grupos (paper, nossa reprodução, WiSARD-family).}
  \label{fig:bar_best_perbase}
\end{figure}

\subsection{Wins matrix: quem lidera em cada base}

A Figura~\ref{fig:wins_matrix} reporta a \emph{wins matrix}: para
cada base, quantos modelos do grupo $X$ ficam dentro de $0{,}02$ do
\emph{best F1 global} entre os três grupos. É uma forma menos
volátil que ``apenas o vencedor'', evitando declarar um grupo
campeão por causa de um \emph{outlier} numericamente irrepetível.

O resultado refina a leitura da média: a liderança concentra-se nas
bases de \emph{leak} saturadas (\texttt{Garage\_Door}, com 8, 7 e 5
modelos no topo nos três grupos; \texttt{Fridge}, com 5 WiSARDs). Nas
bases genuinamente estruturadas (\texttt{Weather}, \texttt{GPS\_Tracker},
\texttt{Modbus}, \texttt{Thermostat}), o topo (dentro de $0{,}02$ do
melhor global) é ocupado só pelos \emph{baselines}: nenhum WiSARD entra,
sinal de que a família chega perto, mas não \emph{iguala}, o melhor
modelo de árvore. A exceção é \texttt{Motion\_Light}, onde apenas a
família pontua, com os dois melhores WiSARDs superando o colapso
majoritário dos \emph{baselines}.

\begin{figure}[ht]
  \centering
  \includegraphics[width=0.85\linewidth]{v06_wins_matrix.png}
  \caption{Wins matrix: número de modelos por grupo a até
           $0{,}02$ F1 do best F1 global da base.}
  \label{fig:wins_matrix}
\end{figure}

\subsection{Distribuição agregada de F1 por modelo}

A Figura~\ref{fig:boxplot_models} é um \emph{boxplot} agregando F1
(média de 4 \emph{folds} por configuração) sobre todas as
configurações e bases do \emph{grid}, por modelo. Modelos com mediana
alta e dispersão baixa são robustos; modelos com cauda longa são
voláteis. RF/CART dominam o topo da distribuição; entre os WiSARDs,
ClusWiSARD tem a mediana mais alta sobre todas as configurações
($0{,}775$), seguida de BTHOWeN ($0{,}739$) e BloomWiSARD ($0{,}733$),
enquanto ULEEN tem a menor mediana ($0{,}603$) e grande dispersão. A
WiSARD pura, apesar do \emph{pico} empatado com ClusWiSARD
(Tabela~\ref{tab:wisard_best_f1}), tem mediana $0{,}701$, a mais
sensível à escolha de hiperparâmetros.

\begin{figure}[ht]
  \centering
  \includegraphics[width=0.85\linewidth]{v06_boxplot_models.png}
  \caption{Distribuição de F1 \emph{weighted} por modelo (todas as
           configurações e bases do \emph{grid}; cada observação é a
           média de 4 \emph{folds} de uma configuração).}
  \label{fig:boxplot_models}
\end{figure}

\subsection{Custo: tempo de \emph{fit} e memória}

As Figuras~\ref{fig:time_per_model} e~\ref{fig:mem} reportam os dois
eixos de custo, medidos sob um \emph{harness} único (\emph{fit} único
por base, incluindo a codificação; \texttt{bench\_cost.py}), o que
torna a comparação \emph{entre} famílias legítima, ao contrário do
tempo agregado do \emph{grid}, que somava quatro \emph{folds} mais o
\emph{held-out}. No eixo de \textbf{tempo} (Figura~\ref{fig:time_per_model}),
o LSTM domina o teto (mediana $\approx 61{,}9$\,s), seguido agora
do ULEEN ($\approx 7{,}0$\,s, mais lento que o SVM em
$\approx 5{,}6$\,s): a configuração \emph{justa} do ULEEN usa cinco
submodelos e vinte \emph{epochs}, e sua implementação de referência é
em \emph{Python} (ver a ressalva de \emph{implementação} adiante). O
trio em C++ treina em $\approx 0{,}45$--$0{,}55$\,s e o BTHOWeN é o mais
rápido da família (mediana $\approx 0{,}09$\,s, cerca de $705\times$
mais rápido que o LSTM), embora RF ($0{,}04$\,s) e kNN ($0{,}004$\,s)
treinem em tempo comparável ou menor.

\begin{figure}[ht]
  \centering
  \includegraphics[width=0.85\linewidth]{v06_time_per_model.png}
  \caption{Tempo mediano de \emph{fit} por modelo (\emph{fit} único,
           mediana sobre 7 bases, escala log). LSTM e ULEEN dominam o
           teto; a família WiSARD fica entre $0{,}09$\,s (BTHOWeN) e
           $7{,}0$\,s (ULEEN).}
  \label{fig:time_per_model}
\end{figure}

No eixo de \textbf{memória} (Figura~\ref{fig:mem}) reportamos uma
\emph{única} métrica de estado implantado (\texttt{deployedSizeBytes}),
comparável entre todos os modelos sem peso: para BloomWiSARD/BTHOWeN/ULEEN
é a tabela de \emph{bits} de tamanho fixo; para WiSARD/ClusWiSARD é o
conjunto mínimo de endereços vistos, \emph{bit-packed} e sem falsos
positivos. Na mediana das sete bases, \emph{toda} a família ocupa
$\approx 0{,}75$--$2{,}5$\,KB (BloomWiSARD $0{,}75$, ULEEN $0{,}94$,
BTHOWeN e WiSARD $\approx 2{,}0$, ClusWiSARD $2{,}5$\,KB), contra
kNN $1{,}1$\,MB, CART $0{,}23$\,MB e RF $3{,}1$\,MB, cerca de três
ordens de magnitude abaixo das árvores. A distinção honesta é
\emph{qualitativa}: as variantes de filtro (Bloom/BTHOWeN/ULEEN) têm
\emph{footprint} \textbf{constante} por construção, enquanto o de
WiSARD/ClusWiSARD é \textbf{dependente dos dados}: $\approx 2$\,KB na
mediana, mas chegando a $\approx 115$\,KB em \texttt{Weather} e
$\approx 230$\,KB em \texttt{Modbus}, onde o endereçamento amplo precisa
memorizar muitos endereços distintos para alcançar F1 $\approx 0{,}93$
(\texttt{Weather}) a $\approx 0{,}95$ (\texttt{Modbus}).
É essa memorização que paga o ganho de F1.

\begin{figure}[ht]
  \centering
  \includegraphics[width=0.85\linewidth]{v06_mem_per_model.png}
  \caption{Memória de estado implantado por modelo
           (\texttt{deployedSizeBytes}, mediana sobre 7 bases, escala
           log). Na mediana a família sem peso ocupa $0{,}75$--$2{,}5$\,KB,
           contra $0{,}2$--$3$\,MB dos \emph{baselines} de árvore; o
           \emph{footprint} de Bloom/BTHOWeN/ULEEN é fixo, o de
           WiSARD/ClusWiSARD é dependente dos dados (até $\approx 230$\,KB
           em \texttt{Modbus}).}
  \label{fig:mem}
\end{figure}

\subsection{Fronteira de Pareto F1 $\times$ custo}

A Figura~\ref{fig:pareto} resume o \emph{trade-off} F1 $\times$ custo em
dois eixos, com cada ponto sendo um modelo na sua melhor configuração
média \emph{per-device}. No eixo de tempo de treino (painel esquerdo),
os modelos de árvore dominam: RF e CART atingem F1 $\approx 0{,}83$ em
$\approx 0{,}04$\,s e o kNN $0{,}789$ em $\approx 0{,}004$\,s; o trio
WiSARD/ClusWiSARD treina em $\approx 0{,}5$\,s com F1 $\approx 0{,}82$ e
o BTHOWeN em $0{,}09$\,s com $0{,}76$; LSTM e ULEEN são dominados
(tempo alto). É no eixo de memória (painel direito) que a família se
distingue, mas como \emph{trade-off}, não como domínio absoluto. As
variantes de filtro (Bloom/BTHOWeN/ULEEN) formam a fronteira de baixo
custo com \emph{footprint} \emph{fixo} de $\approx 0{,}75$--$2$\,KB a
F1 $\approx 0{,}76$, três ordens de magnitude abaixo das árvores;
WiSARD e ClusWiSARD recuperam o melhor F1 da família ($0{,}82$) com
$\approx 2$\,KB \emph{medianos}, ainda muito abaixo das árvores, mas com
\emph{footprint} dependente dos dados que cresce nas bases onde fazem a
memorização. Para \emph{deployment} em \emph{edge} sob restrição
\emph{dura} de memória, as variantes de filtro são a escolha; para
máximo F1 dentro da família a um custo de memória ainda sub-árvore,
WiSARD/ClusWiSARD.

\begin{figure}[ht]
  \centering
  \includegraphics[width=\linewidth]{v06_pareto_wisard.png}
  \caption{Pareto a nível de modelo (melhor configuração média) em dois
           eixos de custo: F1 $\times$ tempo de \emph{fit} (esquerda) e
           F1 $\times$ memória (direita), per-device binário. Os modelos de
           árvore dominam em tempo; a família WiSARD ocupa KBs contra MBs
           dos \emph{baselines} em memória, com \emph{footprint} fixo
           para Bloom/BTHOWeN/ULEEN e dependente dos dados para
           WiSARD/ClusWiSARD. O LSTM é omitido no painel de memória por
           não ter métrica comparável.}
  \label{fig:pareto}
\end{figure}

\subsection{Síntese dos resultados}

A Tabela~\ref{tab:final_consolidada} consolida o resultado em uma
única visão: F1 médio sobre todas as bases per-device, F1
\emph{combined} binário, F1 \emph{combined} multi-classe e tempo
mediano de \emph{fit}.

\begin{table}[ht]
  \centering
  \caption{Síntese final: F1 \emph{weighted} médio (sobre 7 bases $\times$
           4 \emph{folds}) per-device, F1 \emph{combined} binário,
           F1 \emph{combined} multi-classe e tempo mediano de \emph{fit}
           único (\texttt{bench\_cost}, inclui codificação). Negrito =
           melhor por coluna.}
  \label{tab:final_consolidada}
  \small
  \begin{tabular}{lrrrr}
    \toprule
    \textbf{Modelo} & \textbf{F1 per-dev.} & \textbf{F1 comb.\ bin.} & \textbf{F1 comb.\ multi.} & \textbf{Tempo (s, med.)} \\
    \midrule
    \multicolumn{5}{l}{\textit{Baselines}} \\
    LR    & 0{,}620 & 0{,}621 & 0{,}296 & 0{,}01 \\
    LDA   & 0{,}618 & 0{,}629 & 0{,}309 & \textbf{0{,}00} \\
    kNN   & 0{,}789 & 0{,}820 & 0{,}541 & \textbf{0{,}00} \\
    RF    & \textbf{0{,}831} & \textbf{0{,}858} & \textbf{0{,}602} & 0{,}04 \\
    CART  & 0{,}828 & 0{,}854 & 0{,}598 & 0{,}04 \\
    NB    & 0{,}615 & 0{,}601 & 0{,}028 & \textbf{0{,}00} \\
    SVM   & 0{,}659 & 0{,}623 & 0{,}294 & 5{,}64 \\
    LSTM  & 0{,}712 & 0{,}758 & 0{,}449 & 61{,}94 \\
    \midrule
    \multicolumn{5}{l}{\textit{WiSARD-family}} \\
    WiSARD       & 0{,}824 & 0{,}724 & 0{,}413 & 0{,}45 \\
    ClusWiSARD   & 0{,}822 & 0{,}725 & 0{,}413 & 0{,}55 \\
    BloomWiSARD  & 0{,}764 & 0{,}718 & 0{,}405 & 0{,}46 \\
    BTHOWeN      & 0{,}760 & 0{,}514 & 0{,}241 & 0{,}09 \\
    ULEEN        & 0{,}758 & 0{,}745 & 0{,}425 & 7{,}04 \\
    \bottomrule
  \end{tabular}
\end{table}

\subsection{Leitura dos resultados}
\label{sec:leitura}

\paragraph{Per-device binário.} RF/CART dominam ($\approx 0{,}83$),
seguidos por kNN ($0{,}789$); a família WiSARD fica logo atrás, com
WiSARD ($0{,}824$) e ClusWiSARD ($0{,}822$) liderando o grupo
\emph{empatadas}, a cerca de $1$ pp de RF/CART e \emph{à frente} de
kNN, LSTM ($0{,}712$) e SVM ($0{,}659$). Em \texttt{Modbus},
WiSARD/ClusWiSARD ($0{,}95$) chegam a \emph{igualar} RF/CART. Em bases
com \emph{leak} categórico forte (\texttt{Fridge}, \texttt{Garage\_Door}),
todos saturam; em bases sem sinal (\texttt{Motion\_Light}), todos
colapsam para predição majoritária.

\paragraph{Combined binário.} RF/CART (com 10 árvores e Gini)
lideram ($\approx 0{,}85$). A família WiSARD cobre $0{,}51$--$0{,}745$;
o melhor (ULEEN em $0{,}745$) fica $11$ pp atrás de RF e é comparável
ao LSTM ($0{,}758$), com WiSARD, ClusWiSARD e BloomWiSARD agora em
$\approx 0{,}72$ (acima dos $0{,}66$--$0{,}69$ da varredura original);
o BTHOWeN destoa em $0{,}514$.

\paragraph{Combined multi-classe.} O \emph{gap} entre os grupos cresce:
RF/CART em $\approx 0{,}60$; kNN em $0{,}54$; LSTM em $0{,}45$. A
família WiSARD fica entre $0{,}24$ (BTHOWeN) e $0{,}43$ (ULEEN), com
ULEEN ($0{,}425$), WiSARD e ClusWiSARD ($0{,}413$) à frente de
LR/LDA/SVM ($\approx 0{,}30$), bem acima de NB ($0{,}03$), mas ainda
abaixo dos modelos baseados em árvore.

\paragraph{Custo computacional.} A vantagem da família WiSARD não está
no tempo de CPU, e sim na memória. Em tempo de treino, o BTHOWeN é o
mais rápido da família ($0{,}09$\,s mediano por \emph{fit}, cerca de
\textbf{705$\times$ mais rápido que o LSTM}), mas RF ($0{,}04$\,s) e kNN
($0{,}004$\,s) treinam em tempo comparável ou menor com F1 superior, de
modo que a família não domina o eixo de tempo. O ULEEN, com a
configuração \emph{justa} (20 \emph{epochs}, cinco sub-modelos), passa a
gastar $\approx 7$\,s (mais que o SVM), mas isso é em parte
artefato de sua implementação de referência em \emph{Python}, e não do
algoritmo (\textbf{ressalva de implementação}: o trio WiSARD/ClusWiSARD/
BloomWiSARD roda em C++; BTHOWeN e ULEEN em Python, de modo que seu
tempo de CPU superestima o custo de uma implementação otimizada ou em
\emph{hardware} dedicado). É na memória que a separação aparece, sob a
métrica única de estado implantado: as variantes de filtro
(Bloom/BTHOWeN/ULEEN) ocupam $\approx 0{,}75$--$2$\,KB \emph{fixos},
três ordens abaixo das árvores ($0{,}2$--$3$\,MB); WiSARD/ClusWiSARD
recuperam o melhor F1 da família com $\approx 2$\,KB \emph{medianos}
(porém dependentes dos dados). A escolha \emph{dentro} da família é
então um \emph{trade-off}: filtro de \emph{footprint} fixo e mínimo
(Bloom/BTHOWeN/ULEEN, F1 $\approx 0{,}76$) versus o F1 mais alto de
WiSARD/ClusWiSARD ($0{,}82$) a um custo de memória ainda sub-árvore.
